% Source problem: public/notes/物工/問題/2016ap_exam.pdf, PDF page 8.
\documentclass[11pt,a4paper]{article}
\usepackage[margin=24mm]{geometry}
\usepackage{fontspec}
\usepackage{xeCJK}
\setCJKmainfont{HaranoAjiMincho-Regular.otf}
\setCJKsansfont{HaranoAjiGothic-Medium.otf}
\usepackage{amsmath,amssymb}
\usepackage{tikz}
\usetikzlibrary{arrows.meta}
\setlength{\parindent}{1em}
\setlength{\parskip}{0.35em}
\renewcommand{\figurename}{図}
\title{物理工学専攻 院試解答集\\2016年 物理学II 第2問}
\author{政岡凜太郎}
\date{}
\begin{document}
\maketitle

長さ$a$の棒$N$本をそれぞれ右または左に向ける。右向きを正とし、右向きの本数を$n^+$、左向きの本数を$n^-$とすると
\begin{equation}
 n^++n^-=N,\qquad n=n^+-n^-,\qquad X=an.
\end{equation}
従って$n$は$-N,-N+2,\ldots,N$の値をとる。棒の向きを一本ずつ足し合わせた終点が$X$である。
\begin{figure}[htbp]
\centering
\begin{tikzpicture}[>=Stealth,scale=0.85]
 \draw[->] (0,0)--(7.2,0) node[right] {$x$};
 \fill (0,0) circle (2pt) node[below left] {始点 $0$};
 \pgfmathsetmacro{\xsum}{0}
 \pgfmathsetmacro{\ystart}{0.2}
 \draw[dashed,gray] (0,0)--(\xsum,\ystart);
 \foreach \step/\sign in {1/1,2/1,3/-1,4/1,5/-1} {
   \pgfmathsetmacro{\y}{\ystart+0.3*(\step-1)}
   \pgfmathsetmacro{\xnext}{\xsum+2*\sign}
   \ifnum\sign=1
     \draw[->,thick,blue] (\xsum,\y)--(\xnext,\y);
   \else
     \draw[->,thick,red] (\xsum,\y)--(\xnext,\y);
   \fi
   \ifnum\step<5
     \pgfmathsetmacro{\ynext}{\y+0.3}
     \draw[dashed,gray] (\xnext,\y)--(\xnext,\ynext);
   \fi
   \pgfmathsetmacro{\xsum}{\xnext}
 }
 \draw[dashed] (\xsum,1.4)--(\xsum,-0.25);
 \fill (\xsum,0) circle (2pt) node[below right] {終点 $X=a$};
 \node[blue] at (5.0,0.85) {右向き};
 \node[red] at (5.0,1.45) {左向き};
\end{tikzpicture}
\caption{右向き3本、左向き2本の変位を順に足した例。縦方向に並べた各矢印は、左右いずれかの変位 $\pm a$ を表す。}
\end{figure}

\section*{[1] 終点の確率とエントロピー}
向きの列は各棒に二通りあるため、全微視状態数は$2^N$である。終点が$X=na$になるには
\begin{equation}
 n^+=\frac{N+n}{2},\qquad n^-=\frac{N-n}{2}
\end{equation}
でなければならず、$N-n$が偶数である。該当する列の数は、$N$個の場所から右向きの$n^+$個を選ぶ組合せ数
\begin{equation}
 \Omega(n)=\frac{N!}{n^+!n^-!}
\end{equation}
なので
\begin{equation}
 P(n)=\frac{\Omega(n)}{2^N}
 =\frac{1}{2^N}\frac{N!}{[(N+n)/2]![(N-n)/2]!}.
\end{equation}
これは二項分布の全確率であり、許される$n$について$\sum_nP(n)=1$となる。

終点を固定した分子のエントロピーはボルツマンの式$S=k_B\ln\Omega$である。Stirling近似を$n^+,n^-\gg1$で使うと
\begin{align}
 S&\simeq k_B\{N\ln N-n^+\ln n^+-n^-\ln n^-\}\\
 &=k_B\left\{N\ln N-\frac{N+n}{2}\ln\frac{N+n}{2}
                      -\frac{N-n}{2}\ln\frac{N-n}{2}\right\}.
\end{align}
全列数$2^N$は全ての終点を合わせた状態数であり、固定した終点のエントロピーにはその終点を実現する組合せ数$\Omega(n)$を使う。

\section*{[2] 伸長に必要な力}
各微視状態のエネルギーが同じなので、固定された$X$でのヘルムホルツ自由エネルギーは定数から$F=-TS$である。$n^\pm=(N\pm X/a)/2$を用いてエントロピーを$X$で微分すると
\begin{equation}
 \left(\frac{\partial S}{\partial X}\right)_T
 =\frac{k_B}{2a}\ln\frac{n^-}{n^+}.
\end{equation}
従って問題の定義による張力は
\begin{equation}
 \boxed{\tau=\left(\frac{\partial F}{\partial X}\right)_T
 =\frac{k_BT}{2a}\ln\frac{n^+}{n^-}
 =\frac{k_BT}{2a}\ln\frac{Na+X}{Na-X}}.
\end{equation}
$0<X<Na$では$\tau>0$であり、分子を伸ばした状態を保つには右向きに引く必要がある。$X\ll Na$では$\ln[(1+y)/(1-y)]=2y+O(y^3)$より
\begin{equation}
 \boxed{\tau=\frac{k_BT}{a^2N}X+O\!\left(\frac{k_BT X^3}{a^4N^3}\right)}.
\end{equation}
小伸長では有効ばね定数$k_BT/(Na^2)$に比例する復元応答となる。

\section*{[3] 外場中の分配関数と平均値}
右向きと左向きの棒のエネルギーはそれぞれ$-\kappa,+\kappa$である。一本の分配関数は$e^{\beta\kappa}+e^{-\beta\kappa}=2\cosh(\beta\kappa)$なので、独立な$N$本について
\begin{equation}
 Z(\beta)=(2\cosh\beta\kappa)^N,\qquad \beta=\frac1{k_BT}.
\end{equation}
平均エネルギーは$-\partial_\beta\ln Z$から
\begin{equation}
 \langle E\rangle=-N\kappa\tanh(\beta\kappa).
\end{equation}
また各棒のエネルギーと変位には$E_i=-(\kappa/a)\Delta X_i$の関係があるため、$\kappa\ne0$では$\langle X\rangle=-(a/\kappa)\langle E\rangle$であり、直接数えても同じく
\begin{equation}
 \langle X\rangle=Na\tanh(\beta\kappa)
\end{equation}
となる。$|\kappa|\ll k_BT$で$\tanh y=y+O(y^3)$を使うと
\begin{equation}
 \langle E\rangle=-\frac{N\kappa^2}{k_BT}
 +O\!\left(\frac{N\kappa^4}{(k_BT)^3}\right),\qquad
 \langle X\rangle=\frac{Na\kappa}{k_BT}
 +O\!\left(\frac{Na\kappa^3}{(k_BT)^3}\right).
\end{equation}
$\kappa=0$でも両者の極限はゼロである。

\section*{[4] ゆらぎ}
正準分布では$\partial_\beta^2\ln Z=\langle E^2\rangle-\langle E\rangle^2$なので
\begin{equation}
 \boxed{\langle E^2\rangle-\langle E\rangle^2
 =N\kappa^2\operatorname{sech}^2(\beta\kappa)}.
\end{equation}
一本の棒の変位$\Delta X_i$は$\pm a$を取り、平均は$a\tanh(\beta\kappa)$、二乗平均は$a^2$である。従って
\begin{equation}
 \operatorname{Var}(\Delta X_i)
 =a^2-a^2\tanh^2(\beta\kappa)
 =a^2\operatorname{sech}^2(\beta\kappa).
\end{equation}
各棒は独立なので、終点の分散は$N$本の分散を足して
\begin{equation}
 \boxed{\langle X^2\rangle-\langle X\rangle^2
 =Na^2\operatorname{sech}^2(\beta\kappa)}.
\end{equation}
両分散は非負で、$\kappa=0$ではエネルギー分散は0、終点の分散は$Na^2$である。固定した$\kappa\ne0$に対して$T\to0$とすると、両分散はゼロに近づく。
\end{document}
